\section{答案与解析：结构力学基础与技巧班第1次作业答案（几何组成分析）}

\begin{analysisbox}
本节为严格重排版：答案与解析均改写为可编辑 LaTeX 文本，不再保留整页扫描图。
题中体系图可参见后一节“作业题目：几何组成分析”的对应图号。
判断几何组成时，建议先算自由度，再按二元体、两刚片、三刚片和瞬变判据复核；
先找可并入大地的刚片，再检查约束方向是否共线、平行或交于一点，最后判定有无多余约束。
\end{analysisbox}

\subsection*{第 1 页}

\paragraph{1-1}
答案：D。

解析：采用两刚片规则分析。由于 4 根链杆平行且不等长，约束方向退化，体系为几何瞬变体系；同时约束数多于保持几何不变所需约束数，故有两个多余约束。

\paragraph{1-2}
答案：B。

解析：原体系的计算自由度为
\[
W=3\times 3-1\times 2-4=3,
\]
故需要再补充 3 个约束。若在 $A$ 端加入固定支座，则刚片 $ABC$ 先并入大地；随后刚片 $DE$ 与大地由链杆 $BD$ 和铰 $E$ 连接，满足两刚片规则，可并入大地；同理，刚片 $FG$ 与大地也满足两刚片规则。于是原体系成为无多余约束的几何不变体系。

\paragraph{1-3}
答案：C。

解析：刚片 $ABC$ 与大地之间由三根不共线且不交于一点的约束联系，即支座 $A$、链杆 $BD$ 和链杆 $CE$ 连接，满足两刚片规则。因此原体系为无多余联系的几何不变体系。

\paragraph{1-4}
答案：几何常变体系；有 1 个多余约束。

解析：采用两刚片规则分析。刚片 $ABCD$ 与含一个多余约束的刚片 $EFGH$ 之间仅由两根链杆 $CE$ 和 $DF$ 连接，不满足两刚片规则，故整体不能形成几何不变体系。由于局部刚片 $EFGH$ 已含一个多余约束，故原体系为有 1 个多余约束的几何常变体系。

\subsection*{第 2 页}

\paragraph{1-5}
答案：$W=-4$。

解析：按刚片体系计算自由度，
\[
W=5\times 3-6\times 2-1\times 3-4=-4.
\]
负自由度说明约束数超过必要约束数，后续还需结合几何组成规则判断是否为几何不变体系以及多余约束数。

\paragraph{1-6}
答案：$\times$。

解析：可由大地依次增加二元体 $AB$、$ACB$、$ADC$、$DEC$、$DFE$、$EGH$、$EIG$、$FJI$、$FKJ$、$KLF$、$KML$。逐次加入二元体不改变体系的几何不变性质，因此原体系为无多余约束的几何不变体系。

\paragraph{1-7}
解析：刚片 $AD$、$BE$ 可直接并入大地，链杆 $DE$ 为一个多余约束。随后刚片 $EFC$ 与大地由固定端 $C$ 和铰 $E$ 连接，满足两刚片规则且产生两个多余约束。因此原体系为有 3 个多余约束的几何不变体系。

\subsection*{第 3 页}

\paragraph{1-8}
解析：采用三刚片规则分析。刚片 $ABC$ 与刚片 $CDE$ 交于铰 $C$；刚片 $ABC$ 与刚片 $GF$ 由链杆 $BF$、$AG$ 连接，其作用线交于铰 $O_2$；刚片 $CDE$ 与刚片 $FG$ 由链杆 $DF$、$GE$ 连接，其作用线交于铰 $O_1$。三铰 $C,O_1,O_2$ 不共线，满足三刚片规则，可组成一个大刚片。该大刚片再与大地由三根不共线且不交于一点的支座链杆连接，满足两刚片规则。因此原体系为无多余约束的几何不变体系。

\paragraph{1-9}
解析：采用三刚片规则分析。刚片 $ABC$ 与刚片 $DEF$ 由链杆 $BD$、$CF$ 连接，其交点为无穷远铰 $O_1$；刚片 $ABC$ 与大地刚片由支座 $A$、链杆 $CG$ 连接，交于铰 $A$；刚片 $DEF$ 与大地刚片由链杆 $FG$、支座 $E$ 连接，交于铰 $E$。三铰共线，不满足三刚片规则。因此原体系为有 1 个多余约束的几何瞬变体系。

\paragraph{1-10}
解析：计算自由度为
\[
W=12\times 2-25=-1.
\]
采用两刚片规则分析，刚片 $ABCD$ 与刚片 $DEFB$ 由铰 $B$ 和铰 $D$ 连接，相当于四个约束，满足两刚片规则但多出一个约束。因此原体系为有 1 个多余约束的几何不变体系。

\subsection*{第 4 页}

\paragraph{1-11}
解析：先去掉二元体 $AB$、$CEG$。随后，含两个多余约束的刚片 $BCD$ 与大地由三根不共线且不交于一点的链杆 $D$、$CF$ 和 $CG$ 连接，满足两刚片规则。因此原体系为有 2 个多余约束的几何不变体系，即两次超静定。

\paragraph{1-12}
解析：刚片 $ABC$ 与大地由三根不共线且不交于一点的支座链杆 $A$ 和链杆 $CD$ 连接，满足两刚片规则，可并入大地。随后刚片 $FBE$ 与大地由铰 $B$ 和链杆 $F$ 连接，也满足两刚片规则；杆 $GE$ 为多余约束。因此原体系为有 1 个多余约束的几何不变体系。

\paragraph{1-13}
解析：可将刚片 $3456$ 等价为链杆 $34$。随后采用两刚片规则分析，刚片 $14$ 与大地由三根不共线且不交于一点的链杆连接，满足两刚片规则。因此原体系为无多余约束的几何不变体系。

\subsection*{第 5 页}

\paragraph{1-14}
解析：先去掉二元体 $ABC$、$AJ$、$DEF$ 和 $GFH$。随后刚片 $JCDI$ 与大地由三根不共线且不交于一点的支座链杆连接，满足两刚片规则。因此原体系为无多余约束的几何不变体系。

\paragraph{1-15}
解析：计算自由度为
\[
W=7\times 2-13\times 1=1.
\]
采用两刚片规则分析，刚片 $ABC$ 与刚片 $DEF$ 仅由链杆 $CD$、$BE$ 连接，不满足两刚片规则。因此原体系为无多余约束的几何常变体系。

\paragraph{1-16}
解析：先去掉二元体 $ABC$，随后采用三刚片规则分析。刚片 $ADE$ 与大地刚片由链杆 $D$、$EH$ 连接，交于无穷远铰 $O_1$；刚片 $CFG$ 与大地刚片由链杆 $F$、$GH$ 连接，交于铰 $O_2$；刚片 $ADE$ 与刚片 $CFG$ 由链杆 $AC$、$EG$ 连接，交于无穷远铰 $O_3$。三铰共线，不满足三刚片规则。因此原体系为有 1 个多余约束的几何瞬变体系。

\subsection*{第 6 页}

\paragraph{1-17}
解析：计算自由度为
\[
W=4\times 3-3\times 2-7\times 1=-1.
\]
刚片 $AB$ 可先并入大地。随后采用两刚片规则分析，杆 $CD$ 与大地由三根不共线且不交于一点的链杆 $BC$ 和支座 $D$ 连接，满足两刚片规则；杆 $CE$ 为多余约束。因此原体系为有 1 个多余约束的几何不变体系。
